%%%%%%%%%%%%%%%%%%%%%%%%%%%%%%%%%%%%%%%%%%%%%%%%%%%%%%%%%%%%%%%%%%%%%%%%
% Plantilla TFG/TFM
% Universidad de Murcia. Facultad de Informática
% Realizado por: José Manuel Requena Plens
% Modificado: Pablo José Rocamora Zamora
% Contacto: pablojoserocamora@gmail.com
%%%%%%%%%%%%%%%%%%%%%%%%%%%%%%%%%%%%%%%%%%%%%%%%%%%%%%%%%%%%%%%%%%%%%%%%


\chapter{Conclusiones y vías futuras}

Este Trabajo Fin de Grado ha abordado el cuello de botella que el algoritmo \textbf{SAWLC} supone en la generación de tomogramas sintéticos dentro de \textbf{PolNet}, adaptando e implementando \textbf{Tetris} como alternativa basada en correlación cruzada 3D mediante FFT. Sobre el algoritmo original se han añadido los mecanismos necesarios para tener en cuenta las estructuras biológicas preexistentes, muestrear semillas válidas en todo el volumen libre y guiar el crecimiento por clústeres compactos. Adicionalmente se ha desarrollado una versión acelerada por GPU y un banco de pruebas reproducible que compara Tetris CPU, Tetris GPU y SAWLC bajo las mismas condiciones experimentales.\\

La aportación conceptual principal es la \textbf{reformulación de la inserción de proteínas como un problema de optimización global discreta}. Mientras SAWLC plantea la colocación como un proceso estocástico de prueba-y-error cuyo coste crece de forma no lineal con la densidad del volumen, Tetris la formula como la búsqueda del máximo de un mapa de correlación calculado en $O(V \log V)$ por inserción, independientemente del grado de saturación. Los resultados experimentales muestran que Tetris alcanza grados de ocupancia claramente mayores que SAWLC, lo que permite ampliar la escala y la densidad de los conjuntos de datos sintéticos empleados en el entrenamiento de redes neuronales profundas para tareas de segmentación, \textit{particle picking} y \textit{denoising}.\\

La línea de continuación natural y prioritaria del proyecto es la \textbf{integración definitiva de Tetris dentro del \textit{framework} PolNet} como módulo de inserción. Adicionalmente, sería interesante extender el algoritmo más allá de las proteínas solubles para que pueda manejar también otros elementos biológicos del tomograma, como filamentos o proteínas de membrana, lo que requeriría adaptar la construcción del \textit{in-shell template} para incorporar restricciones direccionales y de anclaje a superficie.
